\documentclass[13pt, a4paper]{report}

% ── Packages ──────────────────────────────────────────────────────────────
\usepackage[utf8]{inputenc}
\usepackage[T5]{fontenc}
\usepackage[vietnamese]{babel}
\usepackage{times}              % Times New Roman
\usepackage[
  top=3cm, bottom=3cm,
  left=3.5cm, right=2cm
]{geometry}
\usepackage{setspace}
\doublespacing                  % Dan dong 1.5-2
\usepackage{graphicx}
\usepackage{float}
\usepackage{subcaption}
\usepackage{amsmath, amssymb}
\usepackage{booktabs}
\usepackage{longtable}
\usepackage{hyperref}
\hypersetup{
  colorlinks=true,
  linkcolor=black,
  citecolor=black,
  urlcolor=blue
}
\usepackage{listings}
\usepackage{xcolor}
\usepackage{caption}
\usepackage{indentfirst}
\usepackage{fancyhdr}
\usepackage{titlesec}
\usepackage[numbers]{natbib}

% ── Code listing style ────────────────────────────────────────────────────
\lstset{
  basicstyle=\ttfamily\small,
  breaklines=true,
  frame=single,
  numbers=left,
  numberstyle=\tiny,
  keywordstyle=\color{blue},
  commentstyle=\color{gray},
  stringstyle=\color{orange},
}

% ── Header/Footer ─────────────────────────────────────────────────────────
\pagestyle{fancy}
\fancyhf{}
\fancyhead[R]{\thepage}
\fancyhead[L]{\leftmark}
\renewcommand{\headrulewidth}{0.4pt}

% ── Chapter title format ──────────────────────────────────────────────────
\titleformat{\chapter}[block]
  {\normalfont\Large\bfseries\centering}
  {CHƯƠNG \thechapter:}{0.5em}{\MakeUppercase}
\titlespacing*{\chapter}{0pt}{0pt}{20pt}

\begin{document}

% ══════════════════════════════════════════════════════════════════════════
% TRANG BÌA
% ══════════════════════════════════════════════════════════════════════════
\begin{titlepage}
\begin{center}
  \textbf{ĐẠI HỌC [TÊN TRƯỜNG]}\\
  \textbf{KHOA [TÊN KHOA]}\\[2cm]

  \includegraphics[width=0.25\textwidth]{images/logo.png}\\[1cm]

  {\LARGE\bfseries ĐỒ ÁN MÔN HỌC\\[0.5cm]
  HỌC SÂU / DEEP LEARNING}\\[1.5cm]

  {\Huge\bfseries GestuRhythm v2:\\[0.3cm]
  Hệ Thống Tổng Hợp Âm Nhạc\\
  Tương Tác Thời Gian Thực\\
  Dựa Trên Cử Chỉ Người Dùng}\\[2cm]

  \begin{tabular}{ll}
    \textbf{Giảng viên hướng dẫn:} & [Tên GV] \\[0.3cm]
    \textbf{Sinh viên thực hiện:}  & \\
    \quad [Họ tên SV 1]            & [MSSV 1]  \\
    \quad [Họ tên SV 2]            & [MSSV 2]  \\
    \quad [Họ tên SV 3]            & [MSSV 3]  \\[0.3cm]
    \textbf{Lớp:}                  & [Tên lớp] \\
    \textbf{Học kỳ:}               & [HK], Năm học [XXXX--XXXX] \\
  \end{tabular}\\[2cm]

  \textbf{[Thành phố], tháng [X] năm [XXXX]}
\end{center}
\end{titlepage}

% ══════════════════════════════════════════════════════════════════════════
% LỜI CẢM ƠN
% ══════════════════════════════════════════════════════════════════════════
\chapter*{Lời Cảm Ơn}
\addcontentsline{toc}{chapter}{Lời Cảm Ơn}

Nhóm chúng em xin gửi lời cảm ơn chân thành đến [Tên GV], giảng viên hướng dẫn môn học,
đã tận tình hỗ trợ và định hướng trong suốt quá trình thực hiện đồ án.
Chúng em cũng xin cảm ơn [Tên trường/khoa] đã tạo điều kiện học tập và nghiên cứu.

% ══════════════════════════════════════════════════════════════════════════
% TÓM TẮT
% ══════════════════════════════════════════════════════════════════════════
\chapter*{Tóm Tắt}
\addcontentsline{toc}{chapter}{Tóm Tắt}

Báo cáo này trình bày hệ thống \textbf{GestuRhythm v2} --- một hệ thống tổng hợp
âm nhạc tương tác theo thời gian thực dựa trên phân tích cử chỉ người dùng qua webcam.
Hệ thống sử dụng MediaPipe Holistic để trích xuất chuỗi điểm mốc bàn tay và toàn thân,
sau đó áp dụng mạng Transformer Encoder để học và ánh xạ đặc trưng chuyển động sang
vector cảm xúc (Valence-Arousal). Vector này được dùng để điều kiện hóa (condition)
một mô hình sinh nhạc GPT-style --- đã học ngữ pháp âm nhạc từ 170.000 bản MIDI thật
trong Lakh MIDI Dataset --- thông qua cơ chế Cross-Attention. Kết quả là âm nhạc được
sinh ra vừa phản ánh cảm xúc của người dùng vừa tuân thủ quy luật nhạc lý.

\textbf{Từ khóa:} học sâu, Transformer, sinh nhạc, cảm xúc, cử chỉ, thời gian thực.

% ══════════════════════════════════════════════════════════════════════════
% MỤC LỤC
% ══════════════════════════════════════════════════════════════════════════
\tableofcontents
\listoffigures
\listoftables
\newpage

% ══════════════════════════════════════════════════════════════════════════
% CHƯƠNG 1: GIỚI THIỆU
% ══════════════════════════════════════════════════════════════════════════
\chapter{Giới Thiệu}

\section{Đặt Vấn Đề}

Âm nhạc là ngôn ngữ của cảm xúc. Tuy nhiên, việc tạo ra âm nhạc đòi hỏi
nhiều năm luyện tập kỹ năng nhạc cụ và kiến thức nhạc lý. Câu hỏi đặt ra là:
\textit{Liệu một người không biết nhạc có thể tạo ra âm nhạc có cảm xúc chỉ
bằng cách di chuyển cơ thể?}

Đây chính là động lực của đồ án này.

\section{Mục Tiêu}

\begin{itemize}
  \item Xây dựng hệ thống sinh âm nhạc real-time từ cử chỉ người dùng qua webcam.
  \item Áp dụng kiến trúc Transformer Encoder-Decoder với cơ chế Cross-Attention.
  \item Tự thu thập dataset cặp (cử chỉ -- nhãn cảm xúc) và huấn luyện mô hình từ đầu.
  \item Tích hợp đầu ra vào hệ thống tổng hợp âm thanh MIDI real-time qua FluidSynth.
\end{itemize}

\section{Đóng Góp Chính}

So với các hệ thống tương tự, GestuRhythm v2 có điểm khác biệt:
\begin{enumerate}
  \item \textbf{Tách biệt hai luồng học:} (1) học cảm xúc từ cử chỉ và
        (2) học ngữ pháp âm nhạc từ dữ liệu thật --- sau đó kết hợp qua Cross-Attention.
  \item \textbf{Dataset tự thu thập:} 1853 mẫu từ 2 người với 6 trạng thái cảm xúc.
  \item \textbf{Music Prior từ dữ liệu thật:} huấn luyện trên 5000 bản MIDI từ Lakh MIDI Dataset.
\end{enumerate}

\section{Bố Cục Báo Cáo}

Phần còn lại của báo cáo được tổ chức như sau:
Chương~\ref{chap:related} trình bày các nghiên cứu liên quan.
Chương~\ref{chap:method} mô tả kiến trúc hệ thống và phương pháp.
Chương~\ref{chap:data} trình bày quá trình thu thập và xử lý dữ liệu.
Chương~\ref{chap:exp} trình bày kết quả thực nghiệm.
Chương~\ref{chap:conclude} kết luận và hướng phát triển.

% ══════════════════════════════════════════════════════════════════════════
% CHƯƠNG 2: NGHIÊN CỨU LIÊN QUAN
% ══════════════════════════════════════════════════════════════════════════
\chapter{Nghiên Cứu Liên Quan}
\label{chap:related}

\section{Hệ Thống Điều Khiển Âm Nhạc Bằng Cử Chỉ}

Theremin~\cite{theremin1928} là nhạc cụ điện tử đầu tiên điều khiển bằng cử chỉ
tay trong không khí, ra đời năm 1920. Các nghiên cứu hiện đại như
\cite{nime2022handmate} sử dụng MediaPipe để theo dõi tay trong trình duyệt.

\section{Mô Hình Sinh Nhạc}

Magenta~\cite{magenta2016} của Google DeepMind sử dụng RNN để sinh giai điệu.
Music Transformer~\cite{huang2018music} áp dụng cơ chế Attention để học cấu trúc
âm nhạc dài hạn. Jukebox~\cite{jukebox2020} của OpenAI sinh nhạc thô với giọng hát.

\section{Biểu Diễn Cảm Xúc Trong Âm Nhạc}

Mô hình Valence-Arousal~\cite{russell1980} là khung biểu diễn cảm xúc 2 chiều
phổ biến trong nghiên cứu âm nhạc tình cảm.

% ══════════════════════════════════════════════════════════════════════════
% CHƯƠNG 3: PHƯƠNG PHÁP
% ══════════════════════════════════════════════════════════════════════════
\chapter{Kiến Trúc Hệ Thống}
\label{chap:method}

\section{Tổng Quan Kiến Trúc}

Hệ thống GestuRhythm v2 gồm ba module chính như Hình~\ref{fig:architecture}:

\begin{figure}[H]
  \centering
  \includegraphics[width=0.9\textwidth]{images/architecture.png}
  \caption{Kiến trúc tổng quan của GestuRhythm v2}
  \label{fig:architecture}
\end{figure}

\section{Gesture Emotion Encoder}

Module này nhận chuỗi landmarks bàn tay và cơ thể theo thời gian, trích xuất
vector cảm xúc $\mathbf{e} = [v, a] \in [-1, 1]^2$ trong không gian Valence-Arousal.

\textbf{Kiến trúc:}
\begin{itemize}
  \item Input: $(B, T, 225)$ --- $T=30$ frame, 225 features (21 điểm tay $\times 2$ + 33 điểm pose)
  \item Linear embedding: $225 \rightarrow 128$
  \item Positional embedding: $\text{Embed}(100, 128)$
  \item Transformer Encoder: 3 lớp, 4 attention heads, $d_{ff}=256$, Causal Mask
  \item Output head: $128 \rightarrow 64 \rightarrow 2$, kích hoạt Tanh
\end{itemize}

Causal Mask đảm bảo tại thời điểm $t$, mô hình chỉ sử dụng thông tin từ $[0, t]$,
phù hợp với yêu cầu real-time:
\[
  M_{ij} = \begin{cases} 0 & \text{nếu } j \leq i \\ -\infty & \text{nếu } j > i \end{cases}
\]

\section{Music Prior}

GPT-style autoregressive model học phân phối xác suất nốt nhạc:
\[
  P(\text{note}_t \mid \text{note}_0, \ldots, \text{note}_{t-1})
\]

\textbf{Kiến trúc:}
\begin{itemize}
  \item Vocabulary: 130 token (0--127 = nốt MIDI, 128 = REST, 129 = PAD)
  \item Token embedding: $130 \rightarrow 128$
  \item Transformer Encoder với Causal Mask: 4 lớp, 4 heads
  \item Output: logits $(B, T, 130)$
\end{itemize}

\section{Conditioned Music Decoder}

Cross-Attention kết hợp vector cảm xúc với hidden state của Music Prior:

\[
  \text{Decoder}(\mathbf{e}, \mathbf{h}_{\text{prior}}) =
  \text{TransformerDecoder}(\underbrace{\mathbf{h}_{\text{prior}}}_{\text{Query}},
  \underbrace{W_e \mathbf{e}}_{\text{Key/Value}})
\]

Chỉ có trọng số Cross-Attention được huấn luyện; Encoder và Music Prior bị đóng băng (frozen).

\section{Music Theory Constraints}

Scale mask được áp dụng trước softmax để đảm bảo nốt nhạc không bao giờ lạc điệu:
\[
  \tilde{l}_i = \begin{cases} l_i & \text{nếu nốt } i \in \text{scale} \\ -\infty & \text{ngược lại} \end{cases}
\]

% ══════════════════════════════════════════════════════════════════════════
% CHƯƠNG 4: DỮ LIỆU
% ══════════════════════════════════════════════════════════════════════════
\chapter{Dữ Liệu}
\label{chap:data}

\section{Dataset Cử Chỉ -- Cảm Xúc}

\subsection{Thiết Kế Thu Thập}

Nhóm thiết kế 6 chế độ cảm xúc dựa trên mô hình Valence-Arousal:

\begin{table}[H]
  \centering
  \caption{6 chế độ cảm xúc trong dataset}
  \label{tab:modes}
  \begin{tabular}{clccl}
    \toprule
    \textbf{ID} & \textbf{Tên} & \textbf{Valence} & \textbf{Arousal} & \textbf{Cử chỉ đặc trưng} \\
    \midrule
    0 & HAPPY\_HIGH   & +1.0 & +1.0 & Tay nhanh, rộng, lên cao \\
    1 & HAPPY\_LOW    & +1.0 & -0.5 & Tay chậm, tròn, nhẹ \\
    2 & SAD\_HIGH     & -1.0 & +0.8 & Tay thấp, giật cục, co rút \\
    3 & SAD\_LOW      & -1.0 & -1.0 & Tay thấp nhất, rất chậm \\
    4 & NEUTRAL       &  0.0 &  0.0 & Tay giữa, di chuyển đều \\
    5 & NONE          &  0.0 &  0.0 & Không có tay \\
    \bottomrule
  \end{tabular}
\end{table}

\subsection{Thống Kê Dataset}

\begin{table}[H]
  \centering
  \caption{Thống kê dataset cử chỉ -- cảm xúc}
  \label{tab:dataset}
  \begin{tabular}{lc}
    \toprule
    \textbf{Thuộc tính} & \textbf{Giá trị} \\
    \midrule
    Số người thu thập & 2 (A, C) \\
    Tổng số mẫu & 1853 \\
    Số mẫu mỗi chế độ & $\approx 150$ \\
    Số frame mỗi mẫu & 30 ($\approx$1 giây) \\
    Số features mỗi frame & 225 \\
    \bottomrule
  \end{tabular}
\end{table}

\begin{figure}[H]
  \centering
  \includegraphics[width=\textwidth]{images/dataset_analysis.png}
  \caption{Phân tích phân bố dữ liệu thu thập}
  \label{fig:dataset}
\end{figure}

\section{Lakh MIDI Dataset}

Music Prior được huấn luyện trên 5000 bản nhạc từ Lakh MIDI Clean Dataset~\cite{raffel2016},
bao gồm nhiều thể loại nhạc pop, rock, jazz, classical.

% ══════════════════════════════════════════════════════════════════════════
% CHƯƠNG 5: KẾT QUẢ THỰC NGHIỆM
% ══════════════════════════════════════════════════════════════════════════
\chapter{Kết Quả Thực Nghiệm}
\label{chap:exp}

\section{Cài Đặt Thực Nghiệm}

\begin{table}[H]
  \centering
  \caption{Cài đặt huấn luyện}
  \label{tab:training}
  \begin{tabular}{llll}
    \toprule
    \textbf{Module} & \textbf{Optimizer} & \textbf{Epochs} & \textbf{Batch size} \\
    \midrule
    Emotion Encoder    & AdamW, lr=1e-3 & 80  & 32 \\
    Music Prior        & AdamW, lr=5e-4 & 30  & 64 \\
    Conditioned Decoder & Adam, lr=1e-3 & 50  & 32 \\
    \bottomrule
  \end{tabular}
\end{table}

\section{Kết Quả Gesture Emotion Encoder}

\begin{figure}[H]
  \centering
  \includegraphics[width=\textwidth]{images/training_results.png}
  \caption{Kết quả huấn luyện Gesture Emotion Encoder}
  \label{fig:enc_results}
\end{figure}

\begin{table}[H]
  \centering
  \caption{Đánh giá Gesture Emotion Encoder}
  \label{tab:enc_eval}
  \begin{tabular}{lc}
    \toprule
    \textbf{Metric} & \textbf{Giá trị} \\
    \midrule
    Val Loss (MSE) & 0.0059 \\
    MAE Valence    & [điền sau khi chạy] \\
    MAE Arousal    & [điền sau khi chạy] \\
    \bottomrule
  \end{tabular}
\end{table}

\section{Kết Quả Music Prior}

\begin{figure}[H]
  \centering
  \includegraphics[width=0.9\textwidth]{images/music_prior_results.png}
  \caption{Kết quả huấn luyện Music Prior trên Lakh MIDI}
  \label{fig:prior_results}
\end{figure}

\begin{table}[H]
  \centering
  \caption{Đánh giá Music Prior}
  \label{tab:prior_eval}
  \begin{tabular}{lcc}
    \toprule
    \textbf{Metric} & \textbf{Giá trị} & \textbf{Ý nghĩa} \\
    \midrule
    Val Loss (CrossEntropy) & 1.39 & \\
    Perplexity              & 4.01 & Cần 4 lần đoán để ra đúng nốt \\
    Random baseline         & 4.87 & log(130) \\
    Cải thiện so với random & $\times$32 & \\
    \bottomrule
  \end{tabular}
\end{table}

\section{Kết Quả Conditioned Decoder}

\begin{figure}[H]
  \centering
  \includegraphics[width=\textwidth]{images/conditioned_decoder_results.png}
  \caption{So sánh melody sinh ra với cảm xúc Happy (+0.8,+0.8) vs Sad (-0.8,-0.8)}
  \label{fig:dec_results}
\end{figure}

\section{Demo Hệ Thống Real-time}

\begin{figure}[H]
  \centering
  \includegraphics[width=0.85\textwidth]{images/demo_screenshot.png}
  \caption{Giao diện demo real-time của GestuRhythm v2}
  \label{fig:demo}
\end{figure}

% ══════════════════════════════════════════════════════════════════════════
% CHƯƠNG 6: KẾT LUẬN
% ══════════════════════════════════════════════════════════════════════════
\chapter{Kết Luận và Hướng Phát Triển}
\label{chap:conclude}

\section{Kết Luận}

Đồ án đã xây dựng thành công hệ thống GestuRhythm v2 với các kết quả chính:
\begin{enumerate}
  \item Gesture Emotion Encoder đạt MAE $< 0.1$ trên thang $[-1, 1]$, chứng tỏ
        mô hình học được ánh xạ từ cử chỉ sang không gian cảm xúc.
  \item Music Prior đạt perplexity 4.01, cải thiện $\times$32 so với baseline ngẫu nhiên.
  \item Conditioned Decoder với Cross-Attention tạo ra sự khác biệt có ý nghĩa
        giữa melody sinh ra ở các trạng thái cảm xúc khác nhau.
  \item Hệ thống chạy real-time $< 1$ giây latency trên CPU thông thường.
\end{enumerate}

\section{Hạn Chế}

\begin{itemize}
  \item Dataset nhỏ (1853 mẫu, 2 người) dẫn đến overfitting nhẹ.
  \item Melody sinh ra chưa có cấu trúc câu nhạc rõ ràng.
\end{itemize}

\section{Hướng Phát Triển}

\begin{itemize}
  \item Mở rộng dataset: thêm người thu thập, thêm chế độ cảm xúc.
  \item Triển khai web app trên Hugging Face Spaces.
  \item Thêm tính năng record và export MIDI.
\end{itemize}

% ══════════════════════════════════════════════════════════════════════════
% TÀI LIỆU THAM KHẢO
% ══════════════════════════════════════════════════════════════════════════
\bibliographystyle{plainnat}
\bibliography{references}

\end{document}
